\section{Discussion, limitations, and future work}
\new{The evaluation establishes that HCMO 0.3.0 is consistent, FAIR-scored, and executable: its shapes, competency questions, and round-trip fixture are checked against complete expected results on every change. It does not establish formal alignments with the vocabularies HCMO references, formal ISA conformance, or community adoption; these are stated as open work below rather than implied by the evidence.}

\textbf{Limitations.} HCMO is still an early and evolving resource. Some conceptual boundaries are being stabilised, most notably between raw data, processed measurement, interpreted observation, and biological result, where the same information can be viewed at several levels. The diversity of HCM systems makes a single model that is both broad enough to be reusable and precise enough to describe individual experiments an ongoing balance. \new{Statistics expose this sharply: some systems produce measurements only per cage or group, others per individual animal, so HCMO cannot yet define the experimental unit unambiguously, nor link both cases to ISA in a mutually compatible way.} Alignment to external vocabularies is deliberately partial in this version\new{, and the ISA/OBI/STATO connection is instance-level evidence, not an HCMO class mapping}. The factorial and mixed-model resources are interoperability fixtures, not a general HCMO statistics hierarchy. 

\textbf{Future work.} Planned directions include expert review of definitions and axioms identified by the class and property audits; systematic alignment decisions (reuse as-is vs specialise vs keep HCMO-specific) and bridge modules to adjacent ontologies (e.g.\ OBI, PROV-O, STATO, MEDO) and to taxa/anatomy/device vocabularies; extending the current selective QUDT pattern to additional quantity roles where justified; and populating the ontology with real datasets via the TEATIME network and publishing the corresponding SPARQL queries. HCMO is also designed to underpin tooling that is not a contribution of this paper: \emph{knowledge-graph question answering} with ontology-checked queries \cite{sanou2026,bian2025,butz}; \emph{assisted authoring} through forms generated from the ontology and its shapes; and \emph{vendor-data mapping} of proprietary exports into HCMO with preserved provenance.

\textbf{Conclusion.} \new{HCMO is, to our knowledge, the first ontology to model home-cage monitoring as an integrated system spanning the animal, its housing, its environment, and the technical acquisition chain, delivered as a pinned, validated, and citable resource grounded in the COST TEATIME community. Through community review and the steps above, it aims to become a reference that improves the annotation, comparison, and reuse of HCM data across laboratories.}
